\school{}
\major{}
\student{待填写}{项目小组}
\groupname{葫芦娃除\hspace{-0.18em}BUG}
\thesistitle{SmartSpar BikeHub 重构报告}{智慧共享单车运营调度平台重构实践}
\thesistitleeng{SmartSpar BikeHub Refactoring Report}{Refactoring Practice for a Shared Bicycle Operation and Dispatch Platform}
\thesisadvisor{待填写}
\thesisdate{2026}{6}{15}

\infotype{design}

\infoabstract{%
本文档记录 SmartSpar BikeHub 共享单车运营调度平台的重构过程。本次重构围绕工程可运行性、架构可维护性和演示可验证性展开，优先恢复本地可运行、可构建、可测试的工程基线，并在此基础上完成后端服务层抽取、统一响应与权限装饰器、路径配置集中化、预测模型注册表、预测模块降级启动、前端运营中台化美化、需求管理路由修复以及测试文档沉淀。重构后，前端 lint、生产构建、后端 pytest、登录、需求管理和预测功能均通过验证。本文从重构原因、重构计划、体系结构调整、低层代码改进、运行截图、测试报告、文档修订和小组分工等方面进行说明，为最终提交的源码、文档和演示验收提供依据。
}

\infodrawings{0}
\infowordcount{12000}

\infomaterials{
  \item 重构后源码仓库链接：待填写
  \item 本地运行手册：\code{docs/LOCAL_RUNBOOK.md}
  \item 重构记录：\code{docs/REFACTORING_RECORD.md}
  \item 测试报告：\code{docs/TEST_REPORT.md}
}
